% GENERATED FILE. Edit canonical lessons, then run npm run generate:books.
% canonical-source-hash: fnv1a64:e5adcce7f8cdd896
% canonical-lessons: TA-C427-kalanilai, TA-C427-veppam, TA-C427-veppanilai, TA-C427-koluttu, TA-C427-pulunku

\chapter{Climate and Heat}
\label{ch:ta-climate-and-heat}
\hlchaptermodality{427}

\begin{chapteropening}
\textbf{By the end of this chapter:} \emph{I can talk about climate, heat and temperature, a scorching sun and sultry weather.}

Complete the last lesson of chapter 427: I can talk about climate, heat and temperature, a scorching sun and sultry weather.
\end{chapteropening}

\section[\texorpdfstring{\ta{காலநிலை} (kālanilai)}{காலநிலை (kālanilai)}]{\ta{காலநிலை} (kālanilai) — climate}

\label{lesson:TA-C427-kalanilai}



\begin{quote}
Before the new one: say the Tamil for to fill up, then the Tamil for to gasp.
\end{quote}

\subsection*{You'll want to know: \ta{காலநிலை}}
\textbf{\ta{காலநிலை}} — \emph{kālanilai} — \textquotedblleft{}climate\textquotedblright{}.

\begin{quote}
\ta{இந்த} \ta{ஊரின்} \ta{காலநிலை} \ta{எனக்குப்} \ta{பிடிக்கும்}.
\end{quote}

\begin{quote}
\emph{inta ūriṉ kālanilai eṉakkup piṭikkum.}
\end{quote}

\begin{quote}
I like the climate of this town.
\end{quote}

\ta{காலம்}, a season, + \ta{நிலை}, a state.

\begin{cousinweb}
\textbf{\ta{காலநிலை}} (\emph{kālanilai}) is built from \ta{காலம்}, \textquotedblleft{}time, a season\textquotedblright{}, borrowed from Sanskrit \emph{kāla} (CDIAL 3084), and \ta{நிலை}, \textquotedblleft{}a state, a standing\textquotedblright{} (DEDR 3675).
\end{cousinweb}

\begin{grammarlens}[title={What you've built}]
One more word for getting around.
\end{grammarlens}

\subsection*{Your turn}
\begin{itemize}
\raggedright
  \item \emph{Say it:} \emph{kālanilai}
  \item \emph{Say it:} \emph{kālanilai}, once more
  \item \emph{Say it:} \emph{inta ūriṉ kālanilai eṉakkup piṭikkum.} — I like the climate of this town.
  \item \emph{Recall:} say the Tamil for to fill up, then the Tamil for to gasp, then say \emph{kālanilai} again
\end{itemize}

\begin{tcolorbox}[breakable,colback=teal!4,colframe=teal!35!black,title={Before you move on}]
What does \ta{காலநிலை} mean? (\textquotedblleft{}Climate\textquotedblright{}.) Say it once more.
\end{tcolorbox}
\section[\texorpdfstring{\ta{வெப்பம்} (veppam)}{வெப்பம் (veppam)}]{\ta{வெப்பம்} (veppam) — warmth; heat}

\label{lesson:TA-C427-veppam}



\begin{quote}
Before the new one: say the Tamil for to gasp, then the Tamil for climate.
\end{quote}

\subsection*{You'll want to know: \ta{வெப்பம்}}
\textbf{\ta{வெப்பம்}} — \emph{veppam} — \textquotedblleft{}warmth; heat\textquotedblright{}.

\begin{quote}
\ta{கோடையில்} \ta{வெப்பம்} \ta{அதிகமாக} \ta{இருக்கும்}.
\end{quote}

\begin{quote}
\emph{kōṭaiyil veppam atikamāka irukkum.}
\end{quote}

\begin{quote}
In summer the heat is high.
\end{quote}

A -\ta{அம்} noun: \ta{வெப்பத்தால்}, because of the heat.

\begin{cousinweb}
\textbf{\ta{வெப்பம்}} (\emph{veppam}) is native Dravidian (DEDR 5517), \textquotedblleft{}heat, warmth\textquotedblright{}, from the root \ta{வே}, \textquotedblleft{}to burn, to be hot\textquotedblright{}.
\end{cousinweb}

\begin{grammarlens}[title={What you've built}]
One more word for getting around.
\end{grammarlens}

\subsection*{Your turn}
\begin{itemize}
\raggedright
  \item \emph{Say it:} \emph{veppam}
  \item \emph{Say it:} \emph{veppam}, once more
  \item \emph{Say it:} \emph{kōṭaiyil veppam atikamāka irukkum.} — In summer the heat is high.
  \item \emph{Recall:} say the Tamil for to gasp, then the Tamil for climate, then say \emph{veppam} again
\end{itemize}

\begin{tcolorbox}[breakable,colback=teal!4,colframe=teal!35!black,title={Before you move on}]
What does \ta{வெப்பம்} mean? (\textquotedblleft{}Warmth; heat\textquotedblright{}.) Say it once more.
\end{tcolorbox}
\section[\texorpdfstring{\ta{வெப்பநிலை} (veppanilai)}{வெப்பநிலை (veppanilai)}]{\ta{வெப்பநிலை} (veppanilai) — temperature}

\label{lesson:TA-C427-veppanilai}



\begin{quote}
Before the new one: say the Tamil for climate, then the Tamil for warmth.
\end{quote}

\subsection*{You'll want to know: \ta{வெப்பநிலை}}
\textbf{\ta{வெப்பநிலை}} — \emph{veppanilai} — \textquotedblleft{}temperature\textquotedblright{}.

\begin{quote}
\ta{மதியம்} \ta{வெப்பநிலை} \ta{மிக} \ta{அதிகமாக} \ta{இருந்தது}.
\end{quote}

\begin{quote}
\emph{matiyam veppanilai mika atikamāka iruntatu.}
\end{quote}

\begin{quote}
At noon the temperature was very high.
\end{quote}

\ta{வெப்பம்}, heat, + \ta{நிலை}, a state: the \ta{ம்} drops.

\begin{cousinweb}
\textbf{\ta{வெப்பநிலை}} (\emph{veppanilai}) is built from \ta{வெப்பம்}, \textquotedblleft{}heat\textquotedblright{} (DEDR 5517), and \ta{நிலை}, \textquotedblleft{}a state\textquotedblright{} (DEDR 3675).
\end{cousinweb}

\begin{grammarlens}[title={What you've built}]
One more word for getting around.
\end{grammarlens}

\subsection*{Your turn}
\begin{itemize}
\raggedright
  \item \emph{Say it:} \emph{veppanilai}
  \item \emph{Say it:} \emph{veppanilai}, once more
  \item \emph{Say it:} \emph{matiyam veppanilai mika atikamāka iruntatu.} — At noon the temperature was very high.
  \item \emph{Recall:} say the Tamil for climate, then the Tamil for warmth, then say \emph{veppanilai} again
\end{itemize}

\begin{tcolorbox}[breakable,colback=teal!4,colframe=teal!35!black,title={Before you move on}]
What does \ta{வெப்பநிலை} mean? (\textquotedblleft{}Temperature\textquotedblright{}.) Say it once more.
\end{tcolorbox}
\section[\texorpdfstring{\ta{கொளுத்து} (koḷuttu)}{கொளுத்து (koḷuttu)}]{\ta{கொளுத்து} (koḷuttu) — to scorch; to set alight}

\label{lesson:TA-C427-koluttu}



\begin{quote}
Before the new one: say the Tamil for warmth, then the Tamil for temperature.
\end{quote}

\subsection*{You'll want to know: \ta{கொளுத்து}}
\textbf{\ta{கொளுத்து}} — \emph{koḷuttu} — \textquotedblleft{}to scorch; to set alight\textquotedblright{}.

\begin{quote}
\ta{மதியம்} \ta{வெயில்} \ta{கொளுத்தியது}.
\end{quote}

\begin{quote}
\emph{matiyam veyil koḷuttiyatu.}
\end{quote}

\begin{quote}
At noon the sun was scorching.
\end{quote}

Past (I …): \ta{கொளுத்தினேன்}, \emph{koḷuttiṉēṉ}.

With an object, to light: \ta{விளக்கைக்} \ta{கொளுத்து}, to light a lamp.

\begin{cousinweb}
\textbf{\ta{கொளுத்து}} (\emph{koḷuttu}) is native Dravidian (DEDR 2158), \textquotedblleft{}to kindle, to set on fire\textquotedblright{}.
\end{cousinweb}

\begin{grammarlens}[title={What you've built}]
One more word for getting around.
\end{grammarlens}

\subsection*{Your turn}
\begin{itemize}
\raggedright
  \item \emph{Say it:} \emph{koḷuttu}
  \item \emph{Say it:} \emph{koḷuttu}, once more
  \item \emph{Say it:} \emph{matiyam veyil koḷuttiyatu.} — At noon the sun was scorching.
  \item \emph{Recall:} say the Tamil for warmth, then the Tamil for temperature, then say \emph{koḷuttu} again
\end{itemize}

\begin{tcolorbox}[breakable,colback=teal!4,colframe=teal!35!black,title={Before you move on}]
What does \ta{கொளுத்து} mean? (\textquotedblleft{}To scorch; to set alight\textquotedblright{}.) Say it once more.
\end{tcolorbox}
\section[\texorpdfstring{\ta{புழுங்கு} (puḻuṅku)}{புழுங்கு (puḻuṅku)}]{\ta{புழுங்கு} (puḻuṅku) — to be sultry; to be stuffy}

\label{lesson:TA-C427-pulunku}



\begin{quote}
Before the new one: say the Tamil for temperature, then the Tamil for to scorch.
\end{quote}

\subsection*{You'll want to know: \ta{புழுங்கு}}
\textbf{\ta{புழுங்கு}} — \emph{puḻuṅku} — \textquotedblleft{}to be sultry; to be stuffy\textquotedblright{}.

\begin{quote}
\ta{மழைக்கு} \ta{முன்பு} \ta{மிகவும்} \ta{புழுங்கியது}.
\end{quote}

\begin{quote}
\emph{maḻaikku muṉpu mikavum puḻuṅkiyatu.}
\end{quote}

\begin{quote}
Before the rain it was very sultry.
\end{quote}

Past (it …): \ta{புழுங்கியது}, \emph{puḻuṅkiyatu}.

Of food, to be steamed; \ta{புழுங்கல்} \ta{அரிசி} is parboiled rice.

\begin{cousinweb}
\textbf{\ta{புழுங்கு}} (\emph{puḻuṅku}) is native Dravidian (DEDR 4315), \textquotedblleft{}to be steamed, parboiled; to be sultry\textquotedblright{}.
\end{cousinweb}

\begin{grammarlens}[title={What you've built}]
One more word for getting around.
\end{grammarlens}

\subsection*{Your turn}
\begin{itemize}
\raggedright
  \item \emph{Say it:} \emph{puḻuṅku}
  \item \emph{Say it:} \emph{puḻuṅku}, once more
  \item \emph{Say it:} \emph{maḻaikku muṉpu mikavum puḻuṅkiyatu.} — Before the rain it was very sultry.
  \item \emph{Recall:} say the Tamil for temperature, then the Tamil for to scorch, then say \emph{puḻuṅku} again
\end{itemize}

\begin{tcolorbox}[breakable,colback=teal!4,colframe=teal!35!black,title={Before you move on}]
What does \ta{புழுங்கு} mean? (\textquotedblleft{}To be sultry; to be stuffy\textquotedblright{}.) Say it once more.
\end{tcolorbox}
